\documentclass[11pt]{article}
\usepackage[margin=1in]{geometry}
\usepackage{amsmath,amssymb,amsthm,mathtools,microtype,hyperref}
\hypersetup{colorlinks=true,linkcolor=blue,citecolor=blue,urlcolor=blue,pdftitle={Parity, polynomial degree, and contact-weighted error costs},pdfauthor={Dmitry Vaintrob; Kaarel Hanni; OpenAI Codex}}
\newtheorem{theorem}{Theorem}[section]
\newtheorem{lemma}[theorem]{Lemma}
\newtheorem{corollary}[theorem]{Corollary}
\theoremstyle{remark}\newtheorem{remark}[theorem]{Remark}
\newcommand{\E}{\mathbb E}\newcommand{\Prb}{\mathbb P}
\newcommand{\1}{\mathbf 1}\newcommand{\norm}[2]{\lVert #1\rVert_{#2}}
\newcommand{\dist}{\operatorname{dist}}\newcommand{\sgn}{\operatorname{sign}}
\title{Parity, polynomial degree, and contact-weighted error costs}
\author{Dmitry Vaintrob \and Kaarel H\"anni \and OpenAI Codex}
\date{October 16, 2026}
\begin{document}
\maketitle
\begin{abstract}
We prove a degree lower bound for a polynomial whose sign predicts parity with fixed advantage and whose errors have bounded amplitude cost after weighting by the distance to a coupled contact point. For the mixed quadratic and power cost $\min\{|P(x)|^2/(1+\dist(x,y))^\beta,|P(x)|^{2-\beta}\}$, with $0<\beta\le1/2$, a coupling with Hamming-distance moments of order $\sqrt m$ gives a degree lower bound $m^{\theta_\beta}$. At $\beta=1/2$, the exponent is $\theta_\beta=0.719367\ldots$, without a logarithmic loss. The conclusion is stable under approximation and small exceptional sets of source--contact pairs. We also give a coupling construction for bounded-density subsets of the cube, a logarithmic norm-ratio refinement of the scalar Bernstein inequality, and a relative random-restriction estimate for homogeneous Walsh polynomials. The contact-cost hypothesis is explicit: establishing it for a particular Bayesian posterior or repair problem is a separate geometric and statistical task.
\end{abstract}

\section{The contact-cost theorem}
Let $\mu_m$ be uniform probability measure on $\{-1,1\}^m$. Write
\[
 f(x)=\prod_{i=1}^m x_i,\qquad
 \partial_i P(x)=\frac{P(x)-P(x^{\oplus i})}{2},\qquad
 \Delta=\sum_{i=1}^m\partial_i.
\]
Thus $\Delta\chi_S=|S|\chi_S$. All unqualified norms and expectations use $\mu_m$. For $1<t<2$, put
\[
 s_t=1+\frac{2}{\pi}\arcsin\!\left(\frac2t-1\right).
\]
We use the dimension-free scalar Bernstein inequality
\begin{equation}\label{eq:bernstein}
 \norm{\Delta P}{t}\le C_t u^{s_t}\norm{P}{t}
 \qquad (\deg P\le u).
\end{equation}
This is part of the established theory of polynomial inequalities on the cube \cite{EI}. Section~\ref{sec:log} supplies the needed version from the imaginary-power estimate of Carbonaro and Dragi\v cevi\'c \cite{CD}.

Fix $0<\beta\le1/2$ and put $p=2-\beta$. A contact law is a joint distribution $\pi$ of $(U,V)$ in the cube whose first marginal is exactly $\mu_m$. Set $L=\dist(U,V)$. For a margin-separated function $F$, meaning $|F|\ge1$, define
\[
 E_F=\{x:f(x)F(x)<0\}.
\]
Allow an arbitrary exceptional set $B$ of \emph{pairs} $(U,V)$, and define
\begin{equation}\label{eq:cost}
 \mathcal C_{\pi,B}(F)
 =\E_\pi\!\left[\mathbf 1_{\{U\in E_F,\ (U,V)\notin B\}}
 \min\left\{\frac{|F(U)|^2}{(L+1)^\beta},\ |F(U)|^p\right\}\right].
\end{equation}
In particular, an exceptional set of source points is permitted. No independence between amplitudes, labels, distances, or exceptions is assumed.

\begin{theorem}[Stable contact-cost calibration]\label{thm:main}
Fix $c_0>0$, $K<\infty$, and $C_*<\infty$. Suppose that:
\begin{enumerate}
\item $|F|\ge1$ and $\E[f\sgn F]\ge c_0$;
\item $Q$ has degree at most $u\ge1$ and $\norm{F-Q}{p}\le e\le c_0/2$;
\item $\pi(B)\le\delta$ and
\begin{equation}\label{eq:absorb}
 \delta^{1/p-1/2}(p-1)^{-u/2}\le\frac1{16};
\end{equation}
\item $\mathcal C_{\pi,B}(F)\le K$ and
\begin{equation}\label{eq:distance}
 \norm{L+1}{L^p(\pi)}\le C_*\sqrt m.
\end{equation}
\end{enumerate}
Then
\begin{equation}\label{eq:theta}
 u\ge c\,m^{\theta_\beta},\qquad
 \theta_\beta=\frac{1-\beta/4}{s_p},
\end{equation}
where $c>0$ depends only on the displayed fixed parameters.
\end{theorem}

The exceptional set may depend arbitrarily on $F,Q$ and the contact law. Its small probability is not sufficient on its own; condition~\eqref{eq:absorb} controls its amplitude contribution. For $\delta=\rho^b$, it suffices to have $u\le a\log(1/\rho)$ with $a$ sufficiently small.

\begin{corollary}[Explicit contact exponent]\label{cor:endpoint}
The exponent in Theorem~\ref{thm:main} is
\begin{equation}\label{eq:endpoint}
 \theta_\beta=\frac{1-\beta/4}
 {1+(2/\pi)\arcsin(\beta/(2-\beta))}.
\end{equation}
For $\beta=1/2$, it equals $0.7193671500\ldots$. There is no endpoint logarithmic loss: the single exponent $p=2-\beta$ is used throughout.
\end{corollary}

\subsection{Proof of the calibration theorem}
We first record two simple consequences of parity.

\begin{lemma}[Parity and positivity]\label{lem:parity}
If $|F|\ge1$, $g=\sgn F$, and $\E[fg]\ge c_0$, then
\begin{equation}\label{eq:influence}
 \E[g\Delta F]\ge I(g)\ge mc_0,
 \qquad I(g)=\sum_i\Prb\{g(X)\ne g(X^{\oplus i})\}.
\end{equation}
For every real $Q$ and $1\le t<\infty$,
\begin{equation}\label{eq:positive}
 \norm{Q}{t}\le 2\norm{(fQ)_-}{t}+\frac2m\norm{\Delta Q}{t}.
\end{equation}
\end{lemma}
\begin{proof}
On an edge where $g$ changes sign, margin separation gives
$(g(x)-g(x^{\oplus i}))(F(x)-F(x^{\oplus i}))\ge4$.
Edge symmetrization proves the first inequality in~\eqref{eq:influence}.
Also $\E[fg]=\E[f\partial_i g]$, so
$|\E[fg]|\le\E|\partial_i g|=\operatorname{Inf}_i(g)$ for every $i$.
Summing proves its second inequality.

Let $A$ be normalized cube adjacency and $q=fQ$. Since parity reverses on every edge,
$(I+A)q=2f\Delta Q/m$. Write its right side as $h$. Positivity of $A$ gives
$q_+\le A q_-+|h|$. Contractivity and the triangle inequality give
$\norm{q}{t}\le2\norm{q_-}{t}+\norm{h}{t}$.
\end{proof}

\begin{proof}[Proof of Theorem~\ref{thm:main}]
Set $X=\norm Qp$, and let
\[
 b=|F(U)|\mathbf1_{\{U\in E_F,\ (U,V)\notin B\}},\qquad
 z=\min\{b^2/(L+1)^\beta,b^p\},\qquad W=\norm b{L^p(\pi)}.
\]
Pointwise,
$b\le z^{1/2}(L+1)^{\beta/2}+z^{1/p}$.
Since $2-p=\beta$, H\"older's inequality with exponents $2/p$ and $2/\beta$ yields
\begin{align}
 \norm{z^{1/2}(L+1)^{\beta/2}}{L^p(\pi)}
 &\le(\E_\pi z)^{1/2}\norm{L+1}{L^p(\pi)}^{\beta/2},\label{eq:holder}\\
 W&\le K^{1/2}C_*^{\beta/2}m^{\beta/4}+K^{1/p}
 \le C m^{\beta/4}.\label{eq:W}
\end{align}
Keeping the quadratic part of the contact cost is essential for this bound.

The exact uniform first marginal and H\"older's inequality give
\[
 \norm{Q(U)\mathbf1_B}{L^p(\pi)}
 \le\delta^{1/p-1/2}\norm Q2
 \le\delta^{1/p-1/2}(p-1)^{-u/2}X\le X/16,
\]
where the second inequality is the degree-$u$ hypercontractive comparison.
It follows that
$\norm{F\mathbf1_{E_F}}p\le W+X/16+e$.
The negative-part map is $1$-Lipschitz, so Lemma~\ref{lem:parity} implies
\begin{equation}\label{eq:boot}
 X\le2W+X/8+4e+\frac2m\norm{\Delta Q}p.
\end{equation}
If $2C_pu^{s_p}/m\ge1/4$, then $u\ge c m^{1/s_p}$, which is stronger than the conclusion. Otherwise~\eqref{eq:bernstein} absorbs the last term in~\eqref{eq:boot}, giving $X\le C(W+e)$.

Each $\partial_i$ is an $L^1$ contraction, so
$\norm{\Delta(F-Q)}1\le me$. By~\eqref{eq:influence},
\[
 \norm{\Delta Q}p\ge\E[g\Delta Q]\ge m(c_0-e)\ge mc_0/2.
\]
Combining this with~\eqref{eq:bernstein},~\eqref{eq:W}, and the bound on $X$ proves~\eqref{eq:theta}. Formula~\eqref{eq:endpoint} is the substitution $p=2-\beta$.
\end{proof}

\section{Balanced contacts with short Hamming distance}
The distance assumption has a direct construction when contacts are drawn from a set of fixed positive measure.

\begin{theorem}[A bounded-density coupling]\label{thm:coupling}
If $\nu$ is a probability measure on the $m$-cube with $d\nu/d\mu_m\le C$, there is one coupling $X\sim\mu_m$, $Y\sim\nu$ such that, for every $t\ge0$,
\begin{equation}\label{eq:couplingtail}
 \Prb\{\dist(X,Y)>2\sqrt{m(\log C+2t)}+t\}\le2e^{-t}.
\end{equation}
In particular, for every $q\ge1$,
\[
 \norm{\dist(X,Y)}q\le C_0\sqrt{m(\log C+q)}+C_0q.
\]
\end{theorem}
\begin{proof}
Expose coordinates in a fixed order. For the target measure let
$p_i(y_{<i})=\nu(Y_i=1\mid Y_{<i}=y_{<i})$, $\delta_i=|p_i-1/2|$, and
$S_i=\sum_{j\le i}\delta_j^2$. Let $L_i$ denote the likelihood ratio on the first $i$ coordinates. On a positive-likelihood prefix,
\[
 \E_{\mu_m}[\sqrt{L_i}\mid Y_{<i}]
 =\sqrt{L_{i-1}}\frac{\sqrt{p_i}+\sqrt{1-p_i}}{\sqrt2}.
\]
The square of the last factor is at most $1-\delta_i^2\le e^{-\delta_i^2}$. Thus $\sqrt{L_i}e^{S_i/2}$ is a nonnegative supermartingale under $\mu_m$. Zero-likelihood prefixes may be assigned arbitrary conditional probabilities. Since $L_m\le C$,
\begin{equation}\label{eq:biasenergy}
 \E_\nu e^{S_m/2}
 =\E_{\mu_m}L_m e^{S_m/2}
 \le\sqrt C\,\E_{\mu_m}\sqrt{L_m}e^{S_m/2}\le\sqrt C.
\end{equation}

Given the full joint past, maximally couple a fresh uniform bit $X_i$ with a bit having probability $p_i$ of being $+1$. The conditional probability of disagreement is $\delta_i$. The two resulting marginals are exactly $\mu_m$ and $\nu$. If $D$ is the number of disagreements and $A=\sum_i\delta_i$, the Bernoulli exponential supermartingale gives
$\E e^{D-(e-1)A}\le1$. Since $A\le\sqrt{mS_m}$ and $e-1<2$,
\[
 \Prb\{D>2\sqrt{mS_m}+t\}\le e^{-t}.
\]
Equation~\eqref{eq:biasenergy} gives
$\Prb\{S_m>\log C+2t\}\le e^{-t}$. Their union proves~\eqref{eq:couplingtail}; integration proves the moment estimate.
\end{proof}

Two target laws of bounded density can be coupled by gluing their couplings to a common uniform intermediate point. The triangle inequality then gives the same moment scale. In particular, suppose the two sign classes $S_+,S_-$ of a function have measure at least $2\zeta$, and each contains a contact set $A_s$ of measure at least $\zeta/2$. Couple uniform measure on $S_s$ to uniform measure on $A_s$, and mix with weights $\mu_m(S_s)$. The source marginal is exactly uniform, the target has the same original sign, and its density is at most $2/\zeta$. For fixed $\zeta>0$, the distance condition~\eqref{eq:distance} follows.

\begin{remark}[What the coupling does not prove]
A statement about the two marginals does not control a bad relation depending jointly on source and contact. If a geometric repair theorem was proved under a different contact law, it must be proved again for this coupling, or transferred by a valid uniform statistical argument. Likewise, if $\zeta$ tends to zero, the factor $\sqrt{\log(1/\zeta)}$ in the distance bound must be retained. The main theorem uses fixed contact mass.
\end{remark}

\section{A logarithmic Bernstein refinement}\label{sec:log}
The next estimate makes the norm dependence in scalar Bernstein explicit. Its proof follows the fixed-vector interpolation method of Domelevo et al.\ \cite{Domelevo}, with strip harmonic measure in place of Gaussian damping. We make no priority claim for the interpolation refinement.

\begin{theorem}\label{thm:log}
For $1<t<2$, set $\sigma_t=(2/\pi)\arcsin(2/t-1)$. Every nonzero real cube polynomial of degree at most $u\ge1$ satisfies
\begin{equation}\label{eq:logbern}
 \norm{\Delta P}t\le C_t u
 \left[1+\log\frac{\norm P2}{\norm Pt}\right]^{\sigma_t}\norm Pt.
\end{equation}
The constants are bounded when $t$ ranges over a compact subinterval of $(1,2)$.
\end{theorem}
\begin{proof}
Let $0<\varepsilon\le(t-1)/t$ and define
\[
 q=\frac{t(1-\varepsilon)}{1-t\varepsilon/2},\qquad
 \frac1t=\frac\varepsilon2+\frac{1-\varepsilon}{q}.
\]
We first prove
\begin{equation}\label{eq:mixed}
 \norm{\Delta P}t\le C_tu\varepsilon^{-\sigma_t}
 \norm P2^\varepsilon\norm Pq^{1-\varepsilon}.
\end{equation}
For $L=\Delta+\gamma I$, $\gamma>0$, the imaginary-power estimate of \cite[Proposition 11]{CD} is
\[
 \norm{L^{iv}}{q\to q}\le C(q^*-1)(1+|v|)^{1/2}
 e^{\phi_q|v|},\qquad \phi_q=\arcsin(2/q-1).
\]
Its constants are independent of the dimension and of $\gamma$. Fix a real $g$ with $\norm g{t'}=1$, put $N=1/\varepsilon$, and form the holomorphic bilinear pairing
\[
 G(z)=\E\left[(L^{N(1-z)}P)\,g_z\right],\qquad
 g_z=\sgn(g)|g|^{t'((1-z)/2+z/q')}.
\]
The value is set to zero where $g=0$. At $\theta=1-\varepsilon$, it equals $\E[(LP)g]$. On the left boundary,
$|G(iv)|\le(u+\gamma)^N\norm P2$; on the right,
\[
 |G(1+iv)|\le C(q^*-1)(1+N|v|)^{1/2}
 e^{\phi_qN|v|}\norm Pq.
\]
On the finite cube $G$ is bounded on the closed strip, so the Poisson inequality applies to $\log|G|$.

The right-boundary harmonic-measure density at $\theta$ is
\[
 w_\varepsilon(v)=\frac{\sin(\pi\varepsilon)}{2(\cosh(\pi v)-\cos(\pi\varepsilon))},
 \qquad \int w_\varepsilon=1-\varepsilon.
\]
The elementary bounds needed here are
\begin{align*}
 \int |v|w_\varepsilon(v)\,dv
 &\le\frac{2\varepsilon}{\pi}\bigl(\log(1/\varepsilon)+1\bigr),\\
 \int\log(1+|v|/\varepsilon)w_\varepsilon(v)\,dv&\le C.
\end{align*}
For the first, expand
\[\frac{\sin a}{\cosh x-\cos a}=2\sum_{n\ge1}\sin(na)e^{-nx}\]
and integrate, obtaining
\[(2/\pi^2)\int_0^{\pi\varepsilon}-\log(2\sin(s/2))\,ds.\]
For the second, use $w_\varepsilon(v)\le\pi\varepsilon/[4(v^2+\varepsilon^2)]$ and rescale.

The left boundary has mass $\varepsilon$. Integrating the boundary logarithms consequently gives
\[
 |G(\theta)|\le C_t(u+\gamma)
 \norm P2^\varepsilon\norm Pq^{1-\varepsilon}
 \varepsilon^{-2\phi_q/\pi}.
\]
Here $q\ge1+(t-1)/2$, and $|\phi_q-\phi_t|\le C_t\varepsilon$, so the last power is at most $C_t\varepsilon^{-\sigma_t}$. Taking the dual supremum and then $\gamma\downarrow0$ proves~\eqref{eq:mixed}.

Since $q<t$, $\norm Pq\le\norm Pt$. Take
$\varepsilon=\min\{(t-1)/t,[1+\log(\norm P2/\norm Pt)]^{-1}\}$ to obtain~\eqref{eq:logbern}. All bounds used above are uniform on compact exponent intervals.
\end{proof}

Hypercontractivity gives $\log(\norm P2/\norm Pt)\le(u/2)\log(1/(t-1))$. Substitution recovers~\eqref{eq:bernstein}, since $s_t=1+\sigma_t$. Theorem~\ref{thm:log} therefore does not prove the general linear-degree Bernstein conjecture; it gives linear degree when the norm ratio is bounded.

\section{Homogeneous polynomials under random restrictions}
For completeness we record a relative restriction estimate that avoids any ambient-dimensional norm loss in a structured case. A Walsh polynomial is homogeneous of degree $r$ if all its nonzero Fourier coefficients have cardinality $r$.

\begin{theorem}\label{thm:homogeneous}
Let $1<t\le2$, let $P$ be homogeneous of degree $r$ on the $k$-cube, and choose a uniformly random $m$-coordinate face, with its outside coordinates uniform. Write $P_R$ for the restriction and $(P_R)_{>u}$ for its Fourier part above degree $u$. For $j=u+1$,
\begin{equation}\label{eq:RR}
 \left(\E_R\norm{(P_R)_{>u}}t^t\right)^{1/t}
 \le \binom mj B_t^j\left(\frac rk\right)^{j/2}\norm Pt,
 \qquad B_t=(t-1)^{-1/2}2^{1-t/2}.
\end{equation}
The left side is zero if $j>r$ or $j>m$. The same result holds with arbitrary additional nuisance variables and homogeneity in the active coordinates. A constant plus one homogeneous level obeys the bound with an additional factor $2$.
\end{theorem}
\begin{proof}
Let $\varphi(a)=\sgn(a)|a|^{t-1}$. For real $a,b$,
\[
 (\varphi(a)-\varphi(b))(a-b)
 \ge(t-1)\frac{(a-b)^2}{(|a|+|b|)^{2-t}}.
\]
Edge symmetrization, H\"older's inequality over the input and a uniform coordinate, and $\Delta P=rP$ give
\begin{equation}\label{eq:gradient}
 \frac1k\sum_i\norm{\partial_iP}t^t
 \le B_t^t\left(\frac rk\right)^{t/2}\norm Pt^t.
\end{equation}
For clarity, the symmetrized energy is
$\E[\varphi(P)\Delta P]=(1/4)\sum_i\E[(\varphi(P)-\varphi(P^{\oplus i}))(P-P^{\oplus i})]$.
The other H\"older factor is bounded by
$\E_i\E(|P|+|P^{\oplus i}|)^t\le2^t\norm Pt^t$, which gives exactly the displayed $B_t$.

For distinct derivatives indexed by $S$, the polynomial $\chi_S\partial_SP$ is independent of the coordinates in $S$ and homogeneous of degree $r-|S|$ in the remainder. Iteration of~\eqref{eq:gradient} therefore yields
\[
 \E_{|S|=j}\norm{\partial_SP}t^t
 \le B_t^{jt}\left[\frac{(r)_j}{(k)_j}\right]^{t/2}\norm Pt^t.
\]
Here $(a)_j=a(a-1)\cdots(a-j+1)$.

For any function $F$ on an $m$-cube there is the exact identity
\[
 F_{\ge j}=\sum_{|S|=j}\chi_S\,j\int_0^1(1-z)^{j-1}
 T_z^{[m]\setminus S}(\chi_S\partial_SF)\,dz.
\]
Indeed a Fourier monomial of degree $\ell\ge j$ receives coefficient
$\binom\ell j j\int_0^1z^{\ell-j}(1-z)^{j-1}\,dz=1$.
Each integral operator is an $L^t$ contraction. Minkowski's inequality on the joint space of the random restriction and its input, followed by the preceding derivative bound and $(r)_j/(k)_j\le(r/k)^j$, proves~\eqref{eq:RR}.
For a constant plus one homogeneous level, remove the active conditional mean; the remaining norm is at most twice the original norm.
\end{proof}

\section{Scope and remaining questions}
Theorems~\ref{thm:main} and~\ref{thm:coupling} are independent statements about functions on a finite cube. In applications to kernel or neural-network interpolation, a geometric argument must show that a bounded repair budget controls the cost~\eqref{eq:cost} under the same contact law. If the head or coupling is selected from training data, a uniform statistical argument is also needed. Neither follows from the marginal-density statement alone.

The homogeneous restriction estimate does not extend here to arbitrary sums of Fourier levels. A summable relative restriction bound at effective degree $rm/k$ would imply the general linear Bernstein inequality: average $\Delta_JP=\sum_{u\ge0}(P|_J)_{>u}$ over random sets $J$ of size approximately $k/r$. Thus such an extension is a substantive additional problem. Adaptive normalization on each restricted face can also destroy a direct comparison to the global norm.

The main calibration exponent remains below one. Improving it to a dimension-free obstruction for all degrees $u\le\alpha m$, with a fixed small $\alpha>0$, would require an additional relation between low polynomial degree, amplitude distribution, and contact geometry. This note does not establish a full $\exp(c k\log(ed/k))$ posterior sample lower bound.

\begin{thebibliography}{9}
\bibitem{CD} A. Carbonaro and O. Dragi\v cevi\'c.
\emph{Functional calculus for generators of symmetric contraction semigroups}.
Duke Math. J. 166 (2017), 937--974.
\href{https://arxiv.org/abs/1308.1338}{arXiv:1308.1338}; imaginary-power estimate in Proposition 11 of version 2.
\bibitem{Domelevo} K. Domelevo, P. Durcik, V. Fragkiadaki, O. Klein, D. Oliveira e Silva, L. Slav\'ikov\'a, and B. Wr\'obel.
\emph{Dimension-free estimates for low degree functions on the Hamming cube}.
Studia Math. 280 (2025), 161--173.
\href{https://arxiv.org/abs/2401.07699}{arXiv:2401.07699}.
\bibitem{EI} A. Eskenazis and P. Ivanisvili.
\emph{Polynomial inequalities on the Hamming cube}.
Probab. Theory Related Fields 178 (2020), 235--287.
\href{https://arxiv.org/abs/1902.02406}{arXiv:1902.02406}.
\end{thebibliography}
\end{document}
